\section{Sequential equilibrium at complete-revelation boundaries}
\label{sec:sealed-rounds}

Commit--reveal permits an exact compositional semantics beyond public
owned choice, provided the interface clears the private information.
The values must then retain probabilities. We first give the interface
and its sequential-equilibrium interpretation, then isolate a scalar
fragment and prove that no finite exact winner-observation algebra can
cover even Boolean sealed rounds.

\subsection{A round with one public continuation for every action tuple}

Fix the role set $I$. A round has a finite nonempty action set $A_i$
for every role; inactive roles have a singleton action set. Its syntax is
\[
 \mathsf R_{(A_i)_{i\in I}}(a\mapsto G_a),
 \qquad a\in A=\prod_{i\in I}A_i.
\]
Each role binds its action privately. No action or signal depending on
it becomes available to another role before all actions are bound.
Publication is automatic and reveals the entire tuple to every role.
The game then enters the public continuation $G_a$. Every action tuple,
including one with equilibrium probability zero, has its own continuation.
Leaf games have one winner as before.

There are no private input types, unrevealed cells crossing a round
boundary, common private random seeds, abort choices, or strategic
publication times. Current behavioral randomizations are independent
across roles. The public history can determine the next finite menus
and continuation syntax. Binding and release use the ideal service of
\cref{sec:disclosure}; this is not a theorem about cryptographic or
asynchronous implementation. One active role gives public owned choice.
Two active roles with Boolean menus give a sealed two-bit round.

For commitment and release orders $i_1,\ldots,i_K$ and
$j_1,\ldots,j_K$, the individual round is the disclosure block
\[
 \begin{aligned}
 &\Bind_{i_1}(c_{i_1}:A_{i_1};\{i_1\}).\ \cdots\
       \Bind_{i_K}(c_{i_K}:A_{i_K};\{i_K\}).\\
 &\Pub_{j_1}(c_{j_1};I).\ \cdots\ \Pub_{j_K}(c_{j_K};I).
 \end{aligned}
\]
The selected continuation starts after this block. Audience-derived
views identify exactly the histories with the same preceding public
history at each commitment; no role has a current private input before
its one commitment. Public observations during release are followed by
no further choice until the whole tuple is available. Thus release
order also has no strategic effect in this fragment.

Write $\Delta_I$ for the simplex of winner distributions, and define
\[
 \mathcal E(G)=\{\text{winner distributions of behavioral sequential
 equilibria of }G\}\subseteq\Delta_I.
\]
We use the assessment definition of sequential equilibrium: sequential
rationality together with beliefs that are limits of Bayes beliefs of
completely mixed profiles \cite{KrepsWilson1982}. The coordinates of a
distribution are exactly the players' expected utilities. For a family
of nonempty sets $P_a\subseteq\Delta_I$, define
\begin{equation}\label{se:eq:round}
 \mathcal N_A((P_a)_a)=
 \left\{\sum_{a\in A}\left(\prod_i\sigma_i(a_i)\right)p_a:
 \begin{array}{l}
 p_a\in P_a\text{ for every }a,\\
 \sigma\text{ is a mixed Nash equilibrium of the finite game}\\
 \hfill u_i(a)=(p_a)_i
 \end{array}\right\}.
\end{equation}
Selection of $p_a$ occurs for every tuple, before solving the stage game.
It records the complete continuation strategy, including off-path play;
it is not a choice made by a new strategic agent.

\begin{theorem}[Exact round recursion]\label{se:thm:round}
For the round language,
\[
 \mathcal E(\at i)=\{\delta_i\},\qquad
 \mathcal E(\mathsf R_A(a\mapsto G_a))=
        \mathcal N_A((\mathcal E(G_a))_a).
\]
Each value set is nonempty, compact, and semialgebraic over the rationals.
The round operators are monotone in their continuation sets and commute
with role renaming. Reordering the private commitments within a round
preserves the value set.
\end{theorem}
\begin{proof}
After full publication, each $G_a$ is a proper public subgame. Restricting
an assessment to it gives a sequential equilibrium and a distribution
$p_a$. Any pure change of the current action, followed by the prescribed
continuation, is a legal deviation. The current independent action laws
therefore form a Nash equilibrium for the displayed payoffs.

Conversely choose one continuation equilibrium at every tuple and a
stage Nash profile. Backward induction rules out deviations that also
change later actions: each selected continuation is already optimal
against its opponents, and a current action change cannot improve its
expected continuation payoff. To verify consistency rather than just
subgame optimality, perturb every action law toward its uniform law
with one common positive parameter. At a public round history its reach
factor cancels from Bayes' rule. At a player's commitment, the belief
over preceding hidden commitments is the product of those roles'
perturbed laws. Its limit is the product of the prescribed laws, also
at off-path public histories. Revelation steps offer no choices. This
constructs one consistent assessment for the whole finite game.
The same argument works for every serial order of private commitments.

Nonemptiness follows by choosing continuation distributions and applying
finite mixed Nash existence \cite{Nash1951}. The variables $p_a$ and
$\sigma_i$ range over compact sets. Nash optimality is a finite family
of weak polynomial inequalities, so the feasible set is closed and
its continuous image is compact. The simplex, these inequalities, and
the output equation are rational polynomial conditions. Projection
preserves semialgebraicity, giving the induction. Inclusion of
continuation sets preserves every old feasible choice. Role renaming
only permutes the inequalities and probability coordinates.
\end{proof}

This is the standard backward continuation-payoff construction for games
with observed past actions and simultaneous current moves, expressed as
an algebra for an explicit disclosure boundary. Deterministic open games
provide a related compositional approach \cite{GhaniEtAl2018}.
The proposed mixed extension \cite{GhaniEtAl2020} has finite counterexamples
to its category laws \cite{GameTheoryOpenGameObstructions2026}, detailed
below. It supplies no composition theorem used here. Full revelation
makes the belief calculation local, while the continuation $p_a$ remains
indexed by its public tuple until the stage incentives have been checked.
A retained private cell requires its conditional law in the interface.

\subsection{Finite obstructions to the proposed mixed open-game category}
\label{se:sec:open-game-obstruction}

Ghani, Kupke, Lambert, and Nordvall Forsberg's Definition~11
\cite{GhaniEtAl2020} lifts a history-indexed set of finite strategy laws
$R(x)$ to
\begin{equation}\label{se:eq:relational-lift}
 \mathsf L_R(\mu)=
 \left\{\sum_x\mu(x)\nu_x:
       \nu_x\in\operatorname{conv}_{\!f}R(x)
       \text{ for }x\in\operatorname{supp}\mu\right\}.
\end{equation}
Here $\operatorname{conv}_{\!f}$ is finite convex closure. Their
Definition~12 combines a source strategy law $\mu$ and a later strategy
law $\nu$ by the independent product $\mu\otimes\nu$. The latter must
belong to \eqref{se:eq:relational-lift}, applied to the source's
intermediate-history law and the later equilibrium predicate.
Our examples use identity coplay and a constant zero continuation, so
the first component's payoff averaging imposes no additional condition.
They are valid objects of the proposed open-game structure; their
equilibrium predicates need not be presented as utility maximizers.

\begin{proposition}[Failure of the proposed category laws]
\label{se:prop:open-game-obstruction}
For that raw composition, the left identity law and the canonical
associativity law are false. The two bracketings can have incomparable
equilibrium-law sets. Taking finite convex closure does not yield a
composition congruence, and closing after every composition still does
not restore associativity.
\end{proposition}
\begin{proof}
For the left identity, let a component $K$ have a Boolean strategy,
constant play, and selected laws exactly $\{\delta_0,\delta_1\}$.
At a point history, \eqref{se:eq:relational-lift} gives the finite convex
closure of this set. Consequently $\mathsf{Id};K$ admits the fair law
$(\delta_0+\delta_1)/2$, while $K$ rejects it. This refutes the left-unit
claim in Lemma~16 under the canonical strategy identification.

For associativity put $B=\{0,1\}$. Let $F$ have strategy $s\in B$,
input $\mathsf{Unit}$, output $s$, and every strategy law selected.
Let $H$ have input $b\in B$, strategy $t\in B$, output $t$, and selected
law exactly $\delta_b$. All backward types are $\mathbb R$, coplay
forwards its argument, and selection ignores the continuation.
For a Boolean law $\mu$, write
$d_\mu=\sum_b\mu(b)\delta_{(b,b)}$.
Directly from the two definitions, and identifying the strategy
parentheses canonically, the selected triple laws are
\begin{align*}
 E_{(F;H);H}&=\{\mu\otimes\mu\otimes\mu:\mu\in\mathcal D B\},\\
 E_{F;(H;H)}&=\{\mu\otimes d_\mu:\mu\in\mathcal D B\}.
\end{align*}
Indeed $F;H$ requires independent equal marginals; the second composition
adds a third independent copy. In contrast, $H;H$ at a fixed input $b$
selects only $\delta_{(b,b)}$. Lifting this selection along $\mu$ gives
$d_\mu$, which is then made independent of the source strategy.

Take $u=(\delta_0+\delta_1)/2$. The left bracketing admits the fully
independent fair law $u^{\otimes3}$, assigning $1/8$ to each triple.
The right bracketing rejects it because its last two coordinates agree
almost surely. Conversely $u\otimes d_u$ is admitted on the right and
rejected on the left: its three fair marginals force the sole possible
left candidate $u^{\otimes3}$, which is a different law. This refutes
Theorem~15 and proves incomparability, not merely unequal sets.

Convex equivalence also fails to be a composition congruence. Replace
$F$ by $F_0$, which selects only point laws. It is convex-equivalent
to $F$, but the convex closure of $E_{F_0;H}$ is supported on equal
coordinates, whereas $F;H$ admits $u\otimes u$. For the proposed
repair that closes after every composition, all inputs $F,H$ already
have convex equilibrium sets. The left-associated composite still admits
$u^{\otimes3}$. On the right, $H;H$ at $b$ selects a singleton law;
closing does not change it. Every resulting law, and every finite
mixture of such laws, has its last coordinates equal. It therefore
still excludes $u^{\otimes3}$.
\end{proof}

These are finite counterexamples to specified laws of the published
lifting, not an impossibility theorem for probabilistic open games or
for compositional mixed-game semantics generally. The obstruction is
that a history-indexed choice of strategy laws is averaged into a
marginal and then recombined independently; reassociation changes
which correlations are retained. Restricting to convex equilibrium
sets alone is not closed under the proposed tensor either: the fair
mixture of the point-product laws at $(0,0)$ and $(1,1)$ is correlated.

The pinned \texttt{GameTheory} \texttt{v1-final} artifact
\cite{GameTheoryOpenGameObstructions2026} contains the Lean theorems
\path{probabilistic_leftIdentity_requires_convexity},
\path{probabilistic_seq_not_associative},
\path{probabilistic_seq_bracketings_incomparable},
\path{convexIso_not_left_seq_congruence}, and
\path{saturated_seq_not_associative}. Its \path{Mixed.lean} implements
Definitions~11--12 and records the scope. The independent exact-rational
checker \path{checks/check_open_game_obstruction.py} reproduces the
displayed witnesses. The historical Lean development is not rebuilt
by this paper's \texttt{make check}.

The round recursion in \cref{se:thm:round} uses a different interface:
conditional continuations remain attached to every public action tuple,
and independence is imposed on current simultaneous choices, not on
realized actions across histories. Its SE proof is direct. It asserts
neither a category of arbitrary open components nor the refuted
independent-law composition rules. Its general assessment correspondence
still has the written, rather than kernel-checked, status in
\cref{sec:verification}.

\subsection{Public choice and the limit of winner-set abstraction}

For a distribution set $P$, put
$\operatorname{Win}(P)=\{i:\exists p\in P,\ p_i>0\}$.
The one-active-role instance of \eqref{se:eq:round} is
\begin{equation}\label{se:eq:public}
 P\mathbin{\widehat\star_r}Q=
 \left\{\alpha p+(1-\alpha)q:
 \begin{array}{l}
 p\in P,\ q\in Q,\ 0\leq\alpha\leq1,\\
 \alpha>0\Rightarrow p_r\geq q_r,\\
 \alpha<1\Rightarrow q_r\geq p_r
 \end{array}\right\}.
\end{equation}
For singleton inputs it selects the distribution with larger $r$
coordinate, or the full segment between them on a tie.

\begin{corollary}[Public mixed-support conservativity]\label{se:cor:public}
On finite public owned-choice trees,
$\operatorname{Win}(\mathcal E(G))=B(G)$. Allowing mixed sequential
equilibria adds no possible winner to the pure SPE observation.
\end{corollary}
\begin{proof}
Pure SPE profiles are sequential equilibria in this perfect-information
language. For the other inclusion induct on trees. If a child has pure
winner set $\{r\}$, its mixed continuation distributions are all
$\delta_r$ by induction. The parent owner $r$ can secure utility one,
so every parent equilibrium also has distribution $\delta_r$.
Otherwise the pure recurrence already includes every possible child
winner; mixing can introduce no others. Nonemptiness of each child's
equilibrium set is used throughout.
\end{proof}

Distribution sets need not be convex, even in the public sublanguage.
Let
\[
 T=(\at0\ch2\at1)\ch0\at2.
\]
The left child can give player $0$ any probability $x\in[0,1]$.
When $x>0$, the root must choose left. When $x=0$, the root may mix.
Consequently
\begin{equation}\label{se:eq:nonconvex}
 \mathcal E(T)=
 [\delta_0,\delta_1]\ \cup\ [\delta_1,\delta_2].
\end{equation}
In particular $(\delta_0+\delta_2)/2$ is absent. It belongs to
$\mathcal E(T\ch1 T)$: choose continuation equilibria $\delta_0$ and
$\delta_2$ in the two independent copies, and let the indifferent
root mix equally. Thus distribution-level public choice is not
idempotent. Convexifying value sets would also change the semantics.
The idempotency of the winner-set algebra is a property of its
observation, not of these richer values.

\subsection{A scalar fragment: parallel sum and role exchange}

Fix two active roles $a,b$, with every leaf won by one of them.
These games are constant-sum. In the round language, induction reduces
every stage to a finite zero-sum matrix game. Its equilibrium value
is unique, so write
\[
 \mathcal E(G)=\{v(G)\delta_a+(1-v(G))\delta_b\}.
\]
For public choices owned by $a$ and $b$, the scalar operations are
$\max$ and $\min$. Exchanging the roles throughout a component gives
$C(x)=1-x$.

There is a useful Boolean sealed-round constructor $D(G,H)$:
\begin{center}
\begin{tabular}{c|cc}
 & $b:0$ & $b:1$\\\hline
$a:0$ & $G$ & $\at b$\\
$a:1$ & $\at b$ & $H$
\end{tabular}
\end{center}
The bits are committed privately and both published before entering
the selected continuation. If $v(G)=x$ and $v(H)=y$, its stage payoff
matrix for $a$ is $\left(\begin{smallmatrix}x&0\\0&y\end{smallmatrix}\right)$.

\begin{proposition}[Parallel-sum interpretation]\label{se:prop:parallel}
The constructor $D$ denotes
\[
 x\parallel y=
 \begin{cases}xy/(x+y),&x+y>0,\\0,&x=y=0.\end{cases}
\]
It is associative, commutative, monotone, and zero-absorbing. On the
unit interval it distributes over both $\min$ and $\max$.
\end{proposition}
\begin{proof}
For $x+y>0$, both players assign probability $y/(x+y)$ to their
action $0$. Against this law, either pure action of the opponent gives
$xy/(x+y)$. These two security guarantees force that value in every
Nash equilibrium. If one diagonal entry is zero the same value follows;
if both are zero every strategy profile has value zero.
For positive arguments, reciprocal values turn parallel sum into
addition. This proves associativity and commutativity; the zero cases
are immediate. Monotonicity follows from the displayed formula, and
a monotone unary map on a chain preserves binary minima and maxima.
\end{proof}

Parallel sum is the established operation from network and matrix
theory \cite{AndersonDuffin1969}. Its appearance here follows from a
two-by-two zero-sum game. Its electrical interpretation and algebraic
laws are not new. The constructor has two Boolean commitments and
one occurrence of each continuation; it introduces neither chance nor
numerical rewards at leaves.

\begin{theorem}[Exact rational scalar fragment]\label{se:thm:rational}
The scalar values generated by the winner leaves, role exchange, and
$D$ are exactly $\mathbb Q\cap[0,1]$. In fact $C$ and
$U(x)=x\parallel1=x/(1+x)$ suffice with the leaves.
\end{theorem}
\begin{proof}
The operations preserve rationality. For the converse induct on the
denominator of $p/q$, with $0\leq p\leq q$ and $q>0$. Endpoints are
leaves. If $0<p<q$ and $2p\leq q$, write
\[
 p/q=U\bigl(p/(q-p)\bigr).
\]
Otherwise write
\[
 p/q=C\bigl(U((q-p)/p)\bigr).
\]
Each recursive denominator is smaller. Reduction of fractions is not
needed. This gives an explicit synthesis procedure, without claiming
polynomial size in the binary encoding of $p,q$.
\end{proof}

Rational values of finite rational zero-sum games are standard. The
point here is the small constructor set realizing them with stable
Boolean commitments and complete-revelation boundaries. It supplies
a scalar calculus inside the role-symmetric round semantics, not a
classification of general multiparty equilibrium distributions.

\begin{corollary}[Decidable scalar equations]\label{se:cor:decidable}
Equations in $0,1,C,\parallel,\min,\max$ hold on rational scalar
values exactly when they hold on the real unit interval. Their validity
is decidable by real quantifier elimination.
\end{corollary}
\begin{proof}
All primitives are continuous. Parallel sum is continuous also at
$(0,0)$ because $0\leq x\parallel y\leq\min(x,y)$.
Rational valuations are dense. For the decision claim, the graph of
parallel sum is given on the unit square by $z(x+y)=xy$, with $z=0$
when $x+y=0$; the other graphs are polynomial or piecewise polynomial.
Term evaluation and a failed equality are therefore expressible in the
first-order theory of real closed fields. Its decision theorem applies.
\end{proof}

\subsection{Infinitely many contextual values with three roles}

Put $L_0=\at0$ and $L_{n+1}=D(L_n,\at0)$, using active roles $0,1$.
Then
\[
 v(L_n)=1/(n+1).
\]
Each step adds one round and only one occurrence of its predecessor,
so explicit tree size is linear in $n$. For $n\geq0$ put $G_n=L_{n+1}$.
Every $G_n$ has exactly the possible SE winners $\{0,1\}$ and unique
distribution
\[
 p_n=\frac1{n+2}\delta_0+\frac{n+1}{n+2}\delta_1.
\]
Let $R_m$ be the same construction with roles $0,2$. Consider the
public context $F_m[-]=[-]\ch0 R_m$.

\begin{theorem}[No finite exact winner algebra for sealed rounds]
\label{se:thm:infinite}
For all $n,m\geq0$,
\[
 \operatorname{Win}(\mathcal E(F_m[G_n]))=
 \begin{cases}
 \{0,1\},&n<m,\\
 \{0,1,2\},&n=m,\\
 \{0,2\},&n>m.
 \end{cases}
\]
Thus the possible-winner observation is not compositional on winner
sets. Its contextual equivalence has infinite index, already with
three roles, Boolean commitments, complete revelation after every
round, and no internal role switches. No finite-valued compositional
semantics can preserve even player $2$'s possible-win observation for
this language.
\end{theorem}
\begin{proof}
The root owner compares $1/(n+2)$ with $1/(m+2)$. It selects the
larger value, or may select either continuation on equality. This
gives the formula. All inputs $G_n$ have the same winner set, but the
outputs vary, disproving composition on those sets. In particular
player $2$ is a possible winner exactly when $m\leq n$. Distinct
indices therefore have different responses to the contexts $F_m$.
Any compositional semantics adequate for that observation must assign
them distinct values, so its carrier is infinite. Extra inactive
roles do not affect the construction.
\end{proof}

The first separator uses rates $1/2$ and $1/3$, against an outside
component with probabilities $1/2$ for each of players $0,2$.
The first input ties and permits player $1$ to win; the second loses
the comparison and excludes player $1$. This obstruction does not
depend on subtle off-path beliefs or retained private information.
It already occurs in the fragment where those complications have
been eliminated by full revelation.

\begin{corollary}[Scalar full abstraction]\label{se:cor:scalar}
For closed components in the rational scalar fragment with the same
active pair, equality of scalar values is exactly contextual
equivalence for possible SE winners in the three-role round language.
It suffices to observe the third role's possible victory.
\end{corollary}
\begin{proof}
Equal values give equal singleton distribution sets, which
\eqref{se:eq:round} preserves in every context, including repeated
independent occurrences. For values $v<w$, choose a rational
$v<t<w$. Realize a component of value $t$ between player $a$ and the
third player $c$, using \cref{se:thm:rational}. A public choice owned
by $a$ selects this outside component against value $v$, permitting
$c$ to win, and selects the input against $w$, excluding $c$.
\end{proof}

\subsection{What this interface does and does not compose}

The round operator unifies public and sealed choices in one finite
binder scheme. Full publication allows continuation equilibria to be
selected independently at each public action tuple, and removes a
separate belief parameter from the algebra. Before that boundary,
the commitments cannot be treated as sequential public choices:
revealing one matching-pennies bit before the other's commitment
lets the later player dictate the outcome.

The interface summarizes a closed continuation with the fixed
winner utilities. Replacing a continuation by another with the same
$\mathcal E$ is sound. Replacing the terminal rewards inside a
previously summarized component requires a continuation-dependent
response semantics; it is not supplied by its closed value set.
Likewise a pending secret, private type, selective audience, or
strategic release leaves an information state crossing the interface.
Marginal winner probabilities then do not specify the needed joint
conditional laws and consistent beliefs. A general cut theorem for
that larger language remains open here. The positive recursion and
the infinite-index obstruction identify an exact complete-revelation
fragment, without assigning that fragment the scope of arbitrary
commit--reveal interaction.
